\begin{proof}[Proof of Theorem~\ref{ang:linear-gain}]
We rule out any sequence with subpower input losses, subpower incidence
mass from below, and cover size $\delta^{-q-o(1)}$. This implies the
stated fixed-power result: if no positive $\epsilon,c$ worked, choose
counterexamples with $\epsilon=c=1/j$ at scales
$\delta_j<1/j$ tending to zero. They would form precisely such a
sequence. This observation will also justify all uses of the gain with
sufficiently small fixed losses.

By choosing one of finitely many small projective charts, we can assume
all relevant directions lie in a single chart and that incidence still
has subpower mass. If $q>k-1$, apply
Lemma~\ref{ang:slice-selection} with $h=k-1$. Select a slice with
subpower incidence and normalize its marginal restrictions; their
masses are bounded above by subpower and below by the incidence divided
by the other marginal mass. The coordinate maps take values in a plane, with
\[
 b=q-h\in(0,\tfrac12],\qquad
 f=t-h\in(\tfrac12,1),\qquad b<f,
\]
and tube count at most $\delta^{-b-o(1)}$. Apply
Lemma~\ref{ang:planar-measures} with the original sliced points and
directions as labels, and with coordinate maps $p_{V,z}$ and $v_V$.
Item~3 of Lemma~\ref{ang:slice-selection} gives the ball bounds on their
coordinate pushforwards and the separate cover for each original direction
label. The labeled incidence relation still has subpower mass, so these
data contradict the planar obstruction.

At $q=k-1$, slicing to the plane would leave direction exponent zero,
outside the range of Lemma~\ref{ang:planar-measures}. We therefore put
$h=k-2$, so the selected slices have dimension three, and let
\[
 g=t-h-\tfrac32>0.
\]
Choose once and for all $0<a<g/4$. We will find in one slice three
directions of determinant at least $c\delta^a$ with a common incident
point set of mass at least subpower. The resulting loss in
Lemma~\ref{ang:loomis-whitney} is $O(\delta^{-3a})$, which the
strict gap $g$ will absorb.

Parameterize our chart by $p\in\R^k$, with unit direction
$w(p)=(1,p)/\sqrt{1+|p|^2}$. Restricting and normalizing the chart
costs only a subpower factor; denote the resulting law in $p$ by
$\lambda$. In the rest of the endpoint proof the direction law supplied
to slicing is $w_\#\lambda$, and the incidence relation is restricted to
this chart. The slice law of the chart labels is
\[
 d\lambda_V(p)=\delta^{-h}
 \one_{\{|P_{V^\perp}w(p)|\le\delta\}}\,d\lambda(p),
 \qquad \nu_V=w_\#\lambda_V.
\]
Equation~\eqref{ang:endpoint-plate}
becomes the affine-hyperplane bound in the $p$ coordinates, because
an affine hyperplane there is the intersection of a linear hyperplane
with the chart. The ball exponent remains $k-1$. Choose
\[
 h<\ell<\frac{(k-1)^2}{k}=h+\frac1k.
\]
Apply Lemma~\ref{ang:weak-radial} to this chart law, with a small
fixed output loss $\delta^{-e}$. Its hypotheses hold at sufficiently
small scales for every fixed input-loss threshold. This deletes a
pair set of original product mass $O(\delta^{c_e})$ and gives, for
every anchor $p$, an $\ell$-cap bound $\delta^{-e}r^\ell$ on
the undeleted displacement directions. Also delete pairs at distance
less than $\delta^{b_0}$, where $b_0>0$ will be small. Denote the union
of these two exceptional pair relations by $\mathcal E_{\mathrm{pair}}$,
and call an anchored pair retained when it lies outside this union.

We first record how these deletions transfer to normalized slice-pair
mass. Equation~\eqref{ang:direction-pair-kernel}
shows that for every measurable chart-pair relation $\mathcal E$ with
$(\lambda\otimes\lambda)(\mathcal E)\le\delta^c$, and every fixed
$\theta>0$,
\begin{equation}\label{ang:exceptional-pair-transfer}
 \mathbb E_V(\lambda_V\otimes\lambda_V)(\mathcal E)
 \le C\delta^{c-\theta h}
       +\delta^{\theta(q-h)-o(1)}.
\end{equation}
For $0<\theta\le1$, the first term treats distances greater than
$\delta^\theta$, and the second treats smaller distances by the ball
bound and annular summation. For $\theta\ge1$, the kernel is at most
$C\delta^{-h}$, so the whole expectation is at most
$C\delta^{c-h}\le C\delta^{c-\theta h}$. The same annulus calculation
bounds the expected mass of the close pairs by
$\delta^{b_0(q-h)-o(1)}$. We leave $\theta$ free until $e$, and hence
the weak-radial deletion exponent $c_e$, has been fixed.

We next bound triples that are close to linearly dependent. For two
retained anchored pairs $(p,p'),(p,p'')$, set
\[
 \alpha=\left|
 \frac{p'-p}{|p'-p|}\wedge\frac{p''-p}{|p''-p|}\right|.
\]
In a bounded chart,
\begin{equation}\label{ang:chart-wedge}
 |w(p)\wedge w(p')\wedge w(p'')|
       \asymp |p'-p|\,|p''-p|\,\alpha.
\end{equation}
To see this, put $u=p'-p$ and $v=p''-p$ and subtract the first
unnormalized column from the other two. The orthogonal decomposition
of exterior powers gives the identity
\[
 |(1,p)\wedge(0,u)\wedge(0,v)|^2
       =|u\wedge v|^2+|p\wedge u\wedge v|^2.
\]
It supplies the lower bound, while boundedness of $p$ supplies the
upper bound. Normalizing the three columns changes only fixed constants.
Consequently a wedge smaller than $2\delta^a$ implies
$\alpha\le C\delta^{a-2b_0}$.

For $(w_1,w_2,w_3)=(w(p),w(p'),w(p''))$, the largest singular
value is comparable to one. Their first two columns span area at least $c\delta^{b_0}$,
so the second singular value is at least $c\delta^{b_0}$; by
Equation~\eqref{ang:chart-wedge}, the third is at least
$c\delta^{2b_0}\alpha$. Lemma~\ref{ang:haar} therefore bounds
their normalized expected slicing weight by
\begin{equation}\label{ang:triple-weight}
 \delta^{-3h}\Pr_V\{\max_{1\le i\le3}|P_{V^\perp}w_i|\le\delta\}
 \le C\delta^{-3hb_0}(\alpha+\delta)^{-h}.
\end{equation}
This includes singular values below $\delta$, since the kernel is
truncated there. At fixed $p,p'$, the weak radial estimate and
the two antipodal caps give
\[
 \lambda\{p'':(p,p'')\text{ is retained},\ \alpha\le r\}
       \le C\delta^{-e}r^\ell
       \qquad(\delta\le r\le1).
\]
Let $\mathcal T_\delta$ be the chart triples whose two anchored pairs
are retained and whose original wedge has size less than $2\delta^a$.
Summing dyadic annuli in Equation~\eqref{ang:triple-weight}, up to
$C\delta^{a-2b_0}$, gives
\begin{equation}\label{ang:bad-triples}
 \mathbb E_V\lambda_V^{\otimes3}(\mathcal T_\delta)
 \le C\delta^{(a-2b_0)(\ell-h)-e-3hb_0}.
\end{equation}
Since $\ell>h$, first choose $b_0>0$ sufficiently small depending
on $a,\ell,h$, and then $e>0$ sufficiently small, so that this
exponent is positive. These choices also ensure $a-2b_0>0$.
With $e$ fixed, the weak radial estimate fixes $c_e>0$. Now choose
$\theta>0$ small enough that $c_e-\theta h>0$. Since $q-h=1$, the
second exponent in Equation~\eqref{ang:exceptional-pair-transfer} is
also positive. Together with the close-pair bound, this proves that
$\mathcal E_{\mathrm{pair}}$ has power-small expected slice-pair mass.

Perform the restrictions in Lemma~\ref{ang:slice-selection}, and
write $\widetilde\lambda_V=(w^{-1})_\#\widetilde\nu_V$ for the restricted
law of chart labels. It is dominated by $\lambda_V$ and has the same
uniform subpower mass bound as $\widetilde\nu_V$. Define
\[
 F(V,z,x)=\widetilde\lambda_V
 \{p:(x,w(p))\in\widetilde E_{V,z}\}.
\]
Its integrated first moment is the retained sliced incidence and is at
least subpower, while the total integrated point mass is bounded by a
fixed constant.
H\"older's inequality consequently yields
\begin{equation}\label{ang:triple-incidence}
 \mathbb E_V\int\int F(V,z,x)^3
                 \,d\widetilde\mu_{V,z}(x)\,dz
       \ge\delta^{o(1)}.
\end{equation}
Dropping the point-incidence indicators bounds the contributions of
bad triples and exceptional anchored pairs by the left-hand sides below.
Equation~\eqref{ang:point-slice-total} and
$\widetilde\lambda_V\le\lambda_V$ give
\begin{align*}
 \mathbb E_V\int \|\widetilde\mu_{V,z}\|
      \widetilde\lambda_V^{\otimes3}(\mathcal T_\delta)\,dz
 &\le C\mathbb E_V\lambda_V^{\otimes3}(\mathcal T_\delta),\\
 \mathbb E_V\int \|\widetilde\mu_{V,z}\|\,\|\widetilde\lambda_V\|
      (\widetilde\lambda_V\otimes\widetilde\lambda_V)
          (\mathcal E_{\mathrm{pair}})\,dz
 &\le C\sup_V\|\widetilde\lambda_V\|\,
      \mathbb E_V(\lambda_V\otimes\lambda_V)
          (\mathcal E_{\mathrm{pair}}).
\end{align*}
Both quantities are power-small by Equation~\eqref{ang:bad-triples}
and the preceding bound for $\mathcal E_{\mathrm{pair}}$; the latter
also covers the two possible anchored pairs after changing the constant. No independence
between the point and direction slices is used. Hence a
subpower amount of the triple incidence in
Equation~\eqref{ang:triple-incidence} has original wedge at least
$2\delta^a$. Projection to $V$ and unit normalization change each
vector by $O(\delta)$, so their determinant in $V$ is at least
$\delta^a$ for small $\delta$.

The total parameter mass for integrating $(V,z)$ and three restricted
chart labels is at most subpower, since $\|\widetilde\lambda_V\|$ has a
uniform subpower bound and the $z$-domain is bounded.
It follows that some slice and some such triple have a common
incident point set of sliced mass at least subpower. Let $A$ be its
projected point coordinates. The pushforward of its restricted point
measure has the same mass and inherits the point-coordinate ball bound,
which implies
\[
 |A|_\delta\ge\delta^{-(t-h)+o(1)}.
\]
For each of the three original direction labels, the projected common
point set has its own cover by
$\delta^{-(q-h)-o(1)}=\delta^{-1-o(1)}$ tubes in that label's projected
direction. By
Lemma~\ref{ang:loomis-whitney},
\[
 |A|_\delta\le\delta^{-3/2-3a-o(1)}.
\]
This contradicts $3a<g=t-h-3/2$. When $k=2$ we have $h=0$;
all slicing weights equal one and the same proof applies, with the
$z$-integral interpreted as integration over the singleton
zero-dimensional space. The endpoint and the theorem are proved.
\end{proof}
